\section{What is NES?}
\begin{frame}{NES predicts the next edit, including where it belongs}
\begin{block}{Next Edit Suggestions (NES)}
An IDE continuously uses recent interactions and project context to infer the developer's intent and proactively suggest the next code modification.
\end{block}
\medskip
\small
\begin{tabularx}{\textwidth}{@{}p{3.1cm}XX@{}}
\toprule
 & \textbf{Conventional completion} & \textbf{NES}\\\midrule
Primary signal & Local code near the cursor & Edit stream plus broader project context\\[5pt]
Suggested change & Continuation at the cursor & Insertions, deletions, replacements; potentially other locations or files\\[5pt]
Interaction & Complete the current line & Review, navigate, accept; repeat with Tab\\\bottomrule
\end{tabularx}
\source{Sections 1--2}
\note{Use the paper's distinction as a conceptual comparison. Implementations vary; modern completion tools may also retrieve repository context. NES is distinguished here by proactive next-edit prediction and the repeated per-edit acceptance loop.}
\end{frame}

\begin{frame}{Inside the context-to-edit pipeline}
\centering
\begin{tikzpicture}[
 stage/.style={draw=Teal,fill=Light,rounded corners=2pt,text width=2.05cm,minimum height=1.5cm,align=center,font=\small},
 flow/.style={-{Latex[length=2mm]},thick,Teal}]
\node[stage] (actions) at (0,0) {\textbf{User actions}\\[3pt]Edit, view, move, undo};
\node[stage] (context) at (2.75,0) {\textbf{Context}\\[3pt]History, code, LSP data};
\node[stage] (model) at (5.5,0) {\textbf{NES model}\\[3pt]Predict the next edit};
\node[stage] (diff) at (8.25,0) {\textbf{Parse / diff}\\[3pt]Display the proposed change};
\node[stage] (accept) at (11,0) {\textbf{Developer}\\[3pt]Review, accept, or dismiss};
\draw[flow] (actions) -- (context);
\draw[flow] (context) -- (model);
\draw[flow] (model) -- (diff);
\draw[flow] (diff) -- (accept);
\draw[flow] (accept.south) -- ++(0,-.8) -| (actions.south);
\node[fill=white,text=Teal,font=\small] at (5.5,-1.6) {Accepted edits update the codebase and edit history};
\end{tikzpicture}
\takeaway{The security boundary spans \textbf{triggers, retrieved context, predicted edits, and acceptance}.}
\source{Section 2, Figure 1; simplified architectural redraw}
\note{Walk from the developer's action to event monitoring, context assembly, model prediction, response parsing and Tab acceptance. Acceptance itself becomes another edit in the history, closing the feedback loop.}
\end{frame}
